\subsection{D-instanton Effects in Type IIB String Theory --- arXiv:2205.00609}

\textbf{Authors:} Nathan B. Agmon, Bruno Balthazar, Minjae Cho, Victor A. Rodriguez, Xi Yin

\paragraph{主な主張}
10次元type IIB超弦の散乱振幅について、単一D-instantonの弦結合に関する次次数補正とD/anti-D-instanton対の寄与を求め、S双対性の予言を検証する~\cite{Agmon:2022vdj}。

\paragraph{新規性}
先頭次数を越えるD-instanton振幅を計算し、on-shell計算の曖昧さを超対称性・soft極限および弦の場の理論によって解消する。

\paragraph{位置付け}
D-instanton非摂動効果を世界面計算と弦の場の理論から定量化し、type IIBの双対性と結び付ける研究。

\paragraph{コメント（読書メモ）}
本人の読書メモ：本文冒頭では、弦結合に関する非摂動効果をD-instanton効果と、NSブレーンを含むgravitational-instanton効果の二種類に整理している、と理解した。後者は時空の有効作用に基づくEuclidean運動方程式の解に対応する、場の理論に近いinstanton効果である。

一方、D-instantonは閉弦場だけの作用の鞍点として捉えるのではなく、D-instantonに付着する境界、すなわち「穴」を持つ世界面を通じて散乱振幅へ寄与すると理解した。

したがって、他の非摂動補正と同様に、D-instanton効果を摂動的な結果から区別できる場合に、その寄与を明確に切り分けて議論できる、という点を読書時の要点として記録する。
